% Project 1 — SIMD Advantage Profiling
% Build: `latexmk` in this directory, or `make simd` from the repository root.
\documentclass[11pt]{article}
%% Packages and global settings shared by all ACS reports.
%% Loaded from each project's main.tex:  %% Packages and global settings shared by all ACS reports.
%% Loaded from each project's main.tex:  %% Packages and global settings shared by all ACS reports.
%% Loaded from each project's main.tex:  \input{../common/preamble}

\usepackage[T1]{fontenc}
\usepackage{lmodern}
\usepackage{microtype}
\usepackage[margin=1in]{geometry}
\usepackage{amsmath,amssymb}
\usepackage{graphicx}
\usepackage{booktabs,tabularx,multirow}
\usepackage{siunitx}
\usepackage{subcaption}
\usepackage[dvipsnames]{xcolor}
\usepackage{tcolorbox}
\tcbuselibrary{breakable}
\usepackage{listings}
\usepackage[shortlabels]{enumitem}
\usepackage{csquotes}

% Plots are drawn from the processed CSVs in <project>/data/ (synced from ACS-Code).
\usepackage{pgfplots}
\usepackage{pgfplotstable}
\pgfplotsset{compat=1.18}
\usepgfplotslibrary{groupplots}

\usepackage[backend=bibtex,style=ieee]{biblatex}
\addbibresource{references.bib}

\usepackage[hidelinks]{hyperref}
\usepackage[capitalise,noabbrev]{cleveref}

% ---------------------------------------------------------------- siunitx
\sisetup{
  per-mode = symbol,
  group-digits = integer,
  separate-uncertainty = true,
}
\DeclareSIUnit\flop{FLOP}
\DeclareSIUnit\iops{IOPS}
\DeclareSIUnit\op{op}
\DeclareSIUnit\elem{elem}
\DeclareSIUnit\cycle{cycle}
% Binary prefixes are built in: \qty{128}{\kibi\byte}, \qty{3.2}{\gibi\byte\per\second}

% ---------------------------------------------------------------- listings
\lstset{
  basicstyle = \ttfamily\footnotesize,
  columns = fullflexible,
  keepspaces = true,
  breaklines = true,
  frame = single,
  framerule = 0.3pt,
  rulecolor = \color{black!30},
  showstringspaces = false,
}
% Short disassembly / compiler-report excerpts:  \begin{lstlisting}[style=excerpt] ... \end{lstlisting}
\lstdefinestyle{excerpt}{basicstyle=\ttfamily\scriptsize, frame=leftline}


\usepackage[T1]{fontenc}
\usepackage{lmodern}
\usepackage{microtype}
\usepackage[margin=1in]{geometry}
\usepackage{amsmath,amssymb}
\usepackage{graphicx}
\usepackage{booktabs,tabularx,multirow}
\usepackage{siunitx}
\usepackage{subcaption}
\usepackage[dvipsnames]{xcolor}
\usepackage{tcolorbox}
\tcbuselibrary{breakable}
\usepackage{listings}
\usepackage[shortlabels]{enumitem}
\usepackage{csquotes}

% Plots are drawn from the processed CSVs in <project>/data/ (synced from ACS-Code).
\usepackage{pgfplots}
\usepackage{pgfplotstable}
\pgfplotsset{compat=1.18}
\usepgfplotslibrary{groupplots}

\usepackage[backend=bibtex,style=ieee]{biblatex}
\addbibresource{references.bib}

\usepackage[hidelinks]{hyperref}
\usepackage[capitalise,noabbrev]{cleveref}

% ---------------------------------------------------------------- siunitx
\sisetup{
  per-mode = symbol,
  group-digits = integer,
  separate-uncertainty = true,
}
\DeclareSIUnit\flop{FLOP}
\DeclareSIUnit\iops{IOPS}
\DeclareSIUnit\op{op}
\DeclareSIUnit\elem{elem}
\DeclareSIUnit\cycle{cycle}
% Binary prefixes are built in: \qty{128}{\kibi\byte}, \qty{3.2}{\gibi\byte\per\second}

% ---------------------------------------------------------------- listings
\lstset{
  basicstyle = \ttfamily\footnotesize,
  columns = fullflexible,
  keepspaces = true,
  breaklines = true,
  frame = single,
  framerule = 0.3pt,
  rulecolor = \color{black!30},
  showstringspaces = false,
}
% Short disassembly / compiler-report excerpts:  \begin{lstlisting}[style=excerpt] ... \end{lstlisting}
\lstdefinestyle{excerpt}{basicstyle=\ttfamily\scriptsize, frame=leftline}


\usepackage[T1]{fontenc}
\usepackage{lmodern}
\usepackage{microtype}
\usepackage[margin=1in]{geometry}
\usepackage{amsmath,amssymb}
\usepackage{graphicx}
\usepackage{booktabs,tabularx,multirow}
\usepackage{siunitx}
\usepackage{subcaption}
\usepackage[dvipsnames]{xcolor}
\usepackage{tcolorbox}
\tcbuselibrary{breakable}
\usepackage{listings}
\usepackage[shortlabels]{enumitem}
\usepackage{csquotes}

% Plots are drawn from the processed CSVs in <project>/data/ (synced from ACS-Code).
\usepackage{pgfplots}
\usepackage{pgfplotstable}
\pgfplotsset{compat=1.18}
\usepgfplotslibrary{groupplots}

\usepackage[backend=bibtex,style=ieee]{biblatex}
\addbibresource{references.bib}

\usepackage[hidelinks]{hyperref}
\usepackage[capitalise,noabbrev]{cleveref}

% ---------------------------------------------------------------- siunitx
\sisetup{
  per-mode = symbol,
  group-digits = integer,
  separate-uncertainty = true,
}
\DeclareSIUnit\flop{FLOP}
\DeclareSIUnit\iops{IOPS}
\DeclareSIUnit\op{op}
\DeclareSIUnit\elem{elem}
\DeclareSIUnit\cycle{cycle}
% Binary prefixes are built in: \qty{128}{\kibi\byte}, \qty{3.2}{\gibi\byte\per\second}

% ---------------------------------------------------------------- listings
\lstset{
  basicstyle = \ttfamily\footnotesize,
  columns = fullflexible,
  keepspaces = true,
  breaklines = true,
  frame = single,
  framerule = 0.3pt,
  rulecolor = \color{black!30},
  showstringspaces = false,
}
% Short disassembly / compiler-report excerpts:  \begin{lstlisting}[style=excerpt] ... \end{lstlisting}
\lstdefinestyle{excerpt}{basicstyle=\ttfamily\scriptsize, frame=leftline}

%% Macros shared by all ACS reports.  %% Macros shared by all ACS reports.  %% Macros shared by all ACS reports.  \input{../common/macros} after the preamble.

% ================================================================ draft / final
% Draft (default): rubric boxes are shown and missing data is drawn as a placeholder.
% Final (`make final`): rubric boxes are hidden; TODOs and missing data raise LaTeX warnings.
\newif\ifacsdraft
\acsdrafttrue
\ifdefined\acsfinal \acsdraftfalse \fi

% ================================================================ identity
\newcommand{\acsauthor}{Samson Kempiak}
\newcommand{\acsemail}{kempis@rpi.edu}
\newcommand{\acscourse}{ECSE 4320/6320 Advanced Computer Systems}
\newcommand{\acsdue}{October 21, 2026}
\newcommand{\coderepo}{\href{https://github.com/SKempiak/ACS-Code}{\texttt{github.com/SKempiak/ACS-Code}}}

% ================================================================ writing aids
% \todo{...}   something still to write
% \tbd         a number not measured yet
% \code{...}   file names, commands, flags (underscores are fine)
\newcommand{\todo}[1]{%
  \ifacsdraft\else\PackageWarning{acs}{Unresolved TODO}\fi
  \textcolor{red}{\textbf{[TODO:} #1\textbf{]}}}
\newcommand{\tbd}{%
  \ifacsdraft\else\PackageWarning{acs}{Unmeasured value (\string\tbd)}\fi
  \textcolor{orange}{\textbf{??}}}
\newcommand{\code}[1]{\texttt{\detokenize{#1}}}

% Rubric reminder at the top of a section (draft only):  \rubric{Category}{points}{what is graded}
\newtcolorbox{acsrubricbox}{colback=blue!4, colframe=blue!45, boxrule=0.4pt, arc=1pt,
  left=4pt, right=4pt, top=2pt, bottom=2pt, fontupper=\small}
\newcommand{\rubric}[3]{%
  \ifacsdraft\begin{acsrubricbox}\textbf{Rubric --- #1 (#2 pts).} #3\end{acsrubricbox}\fi}

% Pre-experiment hypothesis (common requirement): write and commit it before measuring.
\newtcolorbox{hypothesis}{breakable, colback=black!3, colframe=black!45, boxrule=0.4pt, arc=1pt,
  left=4pt, right=4pt, top=2pt, bottom=2pt, fonttitle=\bfseries\small,
  title={Hypothesis (from first principles, then compared with the data)}}

% Experiment setup table: one "Knob & value \\" row per line.
%   \begin{knobs}
%     Fixed  & kernel = SAXPY, float32 \\
%     Varied & $N = 2^{8} \ldots 2^{26}$ \\
%   \end{knobs}
\newenvironment{knobs}
  {\par\smallskip\noindent\small\tabularx{\linewidth}{@{}>{\bfseries}l X@{}}\toprule}
  {\bottomrule\endtabularx\par\smallskip}

% ================================================================ data & plots
% Processed CSVs live in <project>/data/<name>.csv (ACS-Code: scripts/sync_data.sh).
%
%   \IfDataExists{baseline}{<tikzpicture using data/baseline.csv>}
%
% draws the plot when the file exists, otherwise a placeholder box.
\newcommand{\acsmissing}[1]{%
  \ifacsdraft\else\PackageWarning{acs}{Missing data file #1}\fi
  \fbox{\parbox[c][0.22\textheight][c]{0.9\linewidth}{\centering\small\color{black!55}%
    Data pending: \texttt{\detokenize{#1}}\\[3pt]
    run the experiment in ACS-Code, then \texttt{scripts/sync\_data.sh}}}}
\newcommand{\IfDataExists}[2]{\IfFileExists{data/#1.csv}{#2}{\acsmissing{data/#1.csv}}}

\pgfplotsset{
  % Base style for every report plot.
  acs/.style={
    width=\linewidth, height=0.55\linewidth,
    grid=major, grid style={black!10},
    legend cell align=left, legend style={font=\footnotesize},
    label style={font=\small}, tick label style={font=\footnotesize},
    unbounded coords=discard,
    table/col sep=comma,
  },
  % Keep only rows where column #1 equals #2. Repeat the key to filter on several columns:
  %   \addplot table[x=n, y=ns_per_elem_median, only if={variant}{autovec}, only if={kernel}{saxpy}] {data/baseline.csv};
  /pgfplots/table/only if/.style 2 args={
    /pgfplots/x filter/.append code={%
      \edef\acsFilterValue{\thisrow{#1}}%
      \edef\acsFilterWant{#2}%
      \ifx\acsFilterValue\acsFilterWant\else\def\pgfmathresult{inf}\fi
    },
  },
  % Error bars spanning the interquartile range of a summarized value (tools/summarize columns):
  %   \addplot+[iqr bars] table[x=n, y=ns_per_elem_median, iqr of=ns_per_elem] {...};
  iqr bars/.style={error bars/.cd, y dir=both, y explicit},
  /pgfplots/table/iqr of/.style={
    y error plus expr={\thisrow{#1_q75}-\thisrow{#1_median}},
    y error minus expr={\thisrow{#1_median}-\thisrow{#1_q25}},
  },
}

% Vertical annotation inside an axis, e.g. cache-level boundaries:  \acsvline{32768}{L1d}
\newcommand{\acsvline}[2]{%
  \draw[densely dashed, black!50] ({axis cs:#1,1}|-{rel axis cs:0,0}) -- ({axis cs:#1,1}|-{rel axis cs:0,1})
    node[pos=0.98, anchor=east, rotate=90, font=\scriptsize, fill=white, inner sep=1pt] {#2};}
 after the preamble.

% ================================================================ draft / final
% Draft (default): rubric boxes are shown and missing data is drawn as a placeholder.
% Final (`make final`): rubric boxes are hidden; TODOs and missing data raise LaTeX warnings.
\newif\ifacsdraft
\acsdrafttrue
\ifdefined\acsfinal \acsdraftfalse \fi

% ================================================================ identity
\newcommand{\acsauthor}{Samson Kempiak}
\newcommand{\acsemail}{kempis@rpi.edu}
\newcommand{\acscourse}{ECSE 4320/6320 Advanced Computer Systems}
\newcommand{\acsdue}{October 21, 2026}
\newcommand{\coderepo}{\href{https://github.com/SKempiak/ACS-Code}{\texttt{github.com/SKempiak/ACS-Code}}}

% ================================================================ writing aids
% \todo{...}   something still to write
% \tbd         a number not measured yet
% \code{...}   file names, commands, flags (underscores are fine)
\newcommand{\todo}[1]{%
  \ifacsdraft\else\PackageWarning{acs}{Unresolved TODO}\fi
  \textcolor{red}{\textbf{[TODO:} #1\textbf{]}}}
\newcommand{\tbd}{%
  \ifacsdraft\else\PackageWarning{acs}{Unmeasured value (\string\tbd)}\fi
  \textcolor{orange}{\textbf{??}}}
\newcommand{\code}[1]{\texttt{\detokenize{#1}}}

% Rubric reminder at the top of a section (draft only):  \rubric{Category}{points}{what is graded}
\newtcolorbox{acsrubricbox}{colback=blue!4, colframe=blue!45, boxrule=0.4pt, arc=1pt,
  left=4pt, right=4pt, top=2pt, bottom=2pt, fontupper=\small}
\newcommand{\rubric}[3]{%
  \ifacsdraft\begin{acsrubricbox}\textbf{Rubric --- #1 (#2 pts).} #3\end{acsrubricbox}\fi}

% Pre-experiment hypothesis (common requirement): write and commit it before measuring.
\newtcolorbox{hypothesis}{breakable, colback=black!3, colframe=black!45, boxrule=0.4pt, arc=1pt,
  left=4pt, right=4pt, top=2pt, bottom=2pt, fonttitle=\bfseries\small,
  title={Hypothesis (from first principles, then compared with the data)}}

% Experiment setup table: one "Knob & value \\" row per line.
%   \begin{knobs}
%     Fixed  & kernel = SAXPY, float32 \\
%     Varied & $N = 2^{8} \ldots 2^{26}$ \\
%   \end{knobs}
\newenvironment{knobs}
  {\par\smallskip\noindent\small\tabularx{\linewidth}{@{}>{\bfseries}l X@{}}\toprule}
  {\bottomrule\endtabularx\par\smallskip}

% ================================================================ data & plots
% Processed CSVs live in <project>/data/<name>.csv (ACS-Code: scripts/sync_data.sh).
%
%   \IfDataExists{baseline}{<tikzpicture using data/baseline.csv>}
%
% draws the plot when the file exists, otherwise a placeholder box.
\newcommand{\acsmissing}[1]{%
  \ifacsdraft\else\PackageWarning{acs}{Missing data file #1}\fi
  \fbox{\parbox[c][0.22\textheight][c]{0.9\linewidth}{\centering\small\color{black!55}%
    Data pending: \texttt{\detokenize{#1}}\\[3pt]
    run the experiment in ACS-Code, then \texttt{scripts/sync\_data.sh}}}}
\newcommand{\IfDataExists}[2]{\IfFileExists{data/#1.csv}{#2}{\acsmissing{data/#1.csv}}}

\pgfplotsset{
  % Base style for every report plot.
  acs/.style={
    width=\linewidth, height=0.55\linewidth,
    grid=major, grid style={black!10},
    legend cell align=left, legend style={font=\footnotesize},
    label style={font=\small}, tick label style={font=\footnotesize},
    unbounded coords=discard,
    table/col sep=comma,
  },
  % Keep only rows where column #1 equals #2. Repeat the key to filter on several columns:
  %   \addplot table[x=n, y=ns_per_elem_median, only if={variant}{autovec}, only if={kernel}{saxpy}] {data/baseline.csv};
  /pgfplots/table/only if/.style 2 args={
    /pgfplots/x filter/.append code={%
      \edef\acsFilterValue{\thisrow{#1}}%
      \edef\acsFilterWant{#2}%
      \ifx\acsFilterValue\acsFilterWant\else\def\pgfmathresult{inf}\fi
    },
  },
  % Error bars spanning the interquartile range of a summarized value (tools/summarize columns):
  %   \addplot+[iqr bars] table[x=n, y=ns_per_elem_median, iqr of=ns_per_elem] {...};
  iqr bars/.style={error bars/.cd, y dir=both, y explicit},
  /pgfplots/table/iqr of/.style={
    y error plus expr={\thisrow{#1_q75}-\thisrow{#1_median}},
    y error minus expr={\thisrow{#1_median}-\thisrow{#1_q25}},
  },
}

% Vertical annotation inside an axis, e.g. cache-level boundaries:  \acsvline{32768}{L1d}
\newcommand{\acsvline}[2]{%
  \draw[densely dashed, black!50] ({axis cs:#1,1}|-{rel axis cs:0,0}) -- ({axis cs:#1,1}|-{rel axis cs:0,1})
    node[pos=0.98, anchor=east, rotate=90, font=\scriptsize, fill=white, inner sep=1pt] {#2};}
 after the preamble.

% ================================================================ draft / final
% Draft (default): rubric boxes are shown and missing data is drawn as a placeholder.
% Final (`make final`): rubric boxes are hidden; TODOs and missing data raise LaTeX warnings.
\newif\ifacsdraft
\acsdrafttrue
\ifdefined\acsfinal \acsdraftfalse \fi

% ================================================================ identity
\newcommand{\acsauthor}{Samson Kempiak}
\newcommand{\acsemail}{kempis@rpi.edu}
\newcommand{\acscourse}{ECSE 4320/6320 Advanced Computer Systems}
\newcommand{\acsdue}{October 21, 2026}
\newcommand{\coderepo}{\href{https://github.com/SKempiak/ACS-Code}{\texttt{github.com/SKempiak/ACS-Code}}}

% ================================================================ writing aids
% \todo{...}   something still to write
% \tbd         a number not measured yet
% \code{...}   file names, commands, flags (underscores are fine)
\newcommand{\todo}[1]{%
  \ifacsdraft\else\PackageWarning{acs}{Unresolved TODO}\fi
  \textcolor{red}{\textbf{[TODO:} #1\textbf{]}}}
\newcommand{\tbd}{%
  \ifacsdraft\else\PackageWarning{acs}{Unmeasured value (\string\tbd)}\fi
  \textcolor{orange}{\textbf{??}}}
\newcommand{\code}[1]{\texttt{\detokenize{#1}}}

% Rubric reminder at the top of a section (draft only):  \rubric{Category}{points}{what is graded}
\newtcolorbox{acsrubricbox}{colback=blue!4, colframe=blue!45, boxrule=0.4pt, arc=1pt,
  left=4pt, right=4pt, top=2pt, bottom=2pt, fontupper=\small}
\newcommand{\rubric}[3]{%
  \ifacsdraft\begin{acsrubricbox}\textbf{Rubric --- #1 (#2 pts).} #3\end{acsrubricbox}\fi}

% Pre-experiment hypothesis (common requirement): write and commit it before measuring.
\newtcolorbox{hypothesis}{breakable, colback=black!3, colframe=black!45, boxrule=0.4pt, arc=1pt,
  left=4pt, right=4pt, top=2pt, bottom=2pt, fonttitle=\bfseries\small,
  title={Hypothesis (from first principles, then compared with the data)}}

% Experiment setup table: one "Knob & value \\" row per line.
%   \begin{knobs}
%     Fixed  & kernel = SAXPY, float32 \\
%     Varied & $N = 2^{8} \ldots 2^{26}$ \\
%   \end{knobs}
\newenvironment{knobs}
  {\par\smallskip\noindent\small\tabularx{\linewidth}{@{}>{\bfseries}l X@{}}\toprule}
  {\bottomrule\endtabularx\par\smallskip}

% ================================================================ data & plots
% Processed CSVs live in <project>/data/<name>.csv (ACS-Code: scripts/sync_data.sh).
%
%   \IfDataExists{baseline}{<tikzpicture using data/baseline.csv>}
%
% draws the plot when the file exists, otherwise a placeholder box.
\newcommand{\acsmissing}[1]{%
  \ifacsdraft\else\PackageWarning{acs}{Missing data file #1}\fi
  \fbox{\parbox[c][0.22\textheight][c]{0.9\linewidth}{\centering\small\color{black!55}%
    Data pending: \texttt{\detokenize{#1}}\\[3pt]
    run the experiment in ACS-Code, then \texttt{scripts/sync\_data.sh}}}}
\newcommand{\IfDataExists}[2]{\IfFileExists{data/#1.csv}{#2}{\acsmissing{data/#1.csv}}}

\pgfplotsset{
  % Base style for every report plot.
  acs/.style={
    width=\linewidth, height=0.55\linewidth,
    grid=major, grid style={black!10},
    legend cell align=left, legend style={font=\footnotesize},
    label style={font=\small}, tick label style={font=\footnotesize},
    unbounded coords=discard,
    table/col sep=comma,
  },
  % Keep only rows where column #1 equals #2. Repeat the key to filter on several columns:
  %   \addplot table[x=n, y=ns_per_elem_median, only if={variant}{autovec}, only if={kernel}{saxpy}] {data/baseline.csv};
  /pgfplots/table/only if/.style 2 args={
    /pgfplots/x filter/.append code={%
      \edef\acsFilterValue{\thisrow{#1}}%
      \edef\acsFilterWant{#2}%
      \ifx\acsFilterValue\acsFilterWant\else\def\pgfmathresult{inf}\fi
    },
  },
  % Error bars spanning the interquartile range of a summarized value (tools/summarize columns):
  %   \addplot+[iqr bars] table[x=n, y=ns_per_elem_median, iqr of=ns_per_elem] {...};
  iqr bars/.style={error bars/.cd, y dir=both, y explicit},
  /pgfplots/table/iqr of/.style={
    y error plus expr={\thisrow{#1_q75}-\thisrow{#1_median}},
    y error minus expr={\thisrow{#1_median}-\thisrow{#1_q25}},
  },
}

% Vertical annotation inside an axis, e.g. cache-level boundaries:  \acsvline{32768}{L1d}
\newcommand{\acsvline}[2]{%
  \draw[densely dashed, black!50] ({axis cs:#1,1}|-{rel axis cs:0,0}) -- ({axis cs:#1,1}|-{rel axis cs:0,1})
    node[pos=0.98, anchor=east, rotate=90, font=\scriptsize, fill=white, inner sep=1pt] {#2};}

%% Measurement environment macros, used by common/environment_table.tex.
%% GENERATED by ACS-Code/scripts/collect_env.sh --tex on the test machine.
%% Regenerate with:
%%   scripts/collect_env.sh --tex > ../ACS-Latex/common/system_info.tex
\newcommand{\sysCPU}{Intel(R) Core(TM) i9-10900F CPU @ 2.80GHz}
\newcommand{\sysArch}{x86\_64}
\newcommand{\sysISA}{sse sse2 ssse3 sse4\_1 sse4\_2 avx avx2 fma f16c -- widest 256-bit (ymm)}
\newcommand{\sysCores}{1 socket(s), 10 cores/socket, 2 threads/core (20 logical)}
\newcommand{\sysSMT}{on (2 threads/core)}
\newcommand{\sysCaches}{L1d 320 KiB (10 instances); L2 2.5 MiB (10 instances); L3 20 MiB (1 instance)}
\newcommand{\sysNUMA}{1 node(s)}
\newcommand{\sysGovernor}{powersave}
\newcommand{\sysDriver}{intel\_pstate}
\newcommand{\sysTurbo}{enabled}
\newcommand{\sysEPP}{balance\_performance}
\newcommand{\sysIsolation}{no isolcpus; runs pinned with taskset/--cpu}
\newcommand{\sysMemory}{63 GiB total}
\newcommand{\sysDIMM}{not available (needs root/dmidecode)}
\newcommand{\sysPageSize}{4096 bytes}
\newcommand{\sysHugePages}{always}
\newcommand{\sysOS}{Arch Linux}
\newcommand{\sysKernel}{7.2.9-arch1-1}
\newcommand{\sysCC}{cc (GCC) 16.2.1 20260810}
\newcommand{\sysGCC}{gcc (GCC) 16.2.1 20260810}
\newcommand{\sysClang}{clang version 23.1.1}
\newcommand{\sysPerf}{not available (not installed)}
\newcommand{\sysPerfParanoid}{2}
\newcommand{\sysFio}{not available (not installed)}
\newcommand{\sysMLC}{not available (not installed)}
\newcommand{\sysGPU}{NVIDIA GeForce RTX 3070}
\newcommand{\sysCUDA}{not available (no CUDA)}
\newcommand{\sysCollected}{2026-10-12 18:35 EDT}


\title{Project 1: SIMD Advantage Profiling}
\author{\acsauthor\\[2pt]\small\acscourse\\\small\href{mailto:\acsemail}{\acsemail}}
\date{\acsdue}

\begin{document}
\maketitle

\begin{abstract}
  We quantify scalar-versus-SIMD speedup for four single-threaded numeric kernels (SAXPY, dot
  product, elementwise multiply and a 3-point stencil) on an Intel Core i9-10900F (Comet Lake,
  AVX2+FMA, \qty{256}{bit} vectors). With cache-resident data a \code{-march=native} build reaches
  about $5.8\times$ over scalar for SAXPY (\num{8} float32 lanes), while the generic
  auto-vectorized (SSE, \qty{128}{bit}) build reaches about $3\times$; these gains compress to
  $\approx1.3\text{--}1.4\times$ once the working set spills to DRAM and memory bandwidth
  dominates. Float64 roughly halves the SIMD advantage (\num{4} lanes), and non-unit stride or
  gather access erases it entirely. Our negative case is memory-bound SAXPY: the loop vectorizes
  (verified by disassembly) yet barely speeds up; a second, subtler anomaly is the dot product,
  whose reduction the compiler refuses to reassociate without \code{-ffast-math}, leaving a scalar
  dependency chain and $\approx1\times$ speedup despite packed multiplies. A roofline built from our
  own measured streaming bandwidth ($\approx\qty{43}{\giga\byte\per\second}$) explains every kernel
  as memory-bound in DRAM.
\end{abstract}

\ifacsdraft\tableofcontents\clearpage\fi

% !TeX root = ../main.tex
\section{Introduction}\label{sec:intro}

Single-instruction-multiple-data (SIMD) units let one instruction operate on a whole vector of
elements. On the Intel Core i9-10900F studied here the widest available unit is AVX2 with FMA, a
\qty{256}{bit} register that holds \num{8} float32 or \num{4} float64 lanes, so the \emph{ideal}
speedup over scalar code is $8\times$ (float32) or $4\times$ (float64). Real kernels rarely reach
that ceiling: loads and stores, loop prologue/epilogue handling, reduction dependencies, and---above
all---finite memory bandwidth cap the benefit. The purpose of this report is to measure where the
benefit appears, by how much it falls short of the lane count, and exactly when it disappears.

This report answers:
\begin{enumerate}
  \item How much faster are auto-vectorized builds than scalar builds for each kernel, and how does
        that change from cache-resident to DRAM-resident working sets? (\cref{sec:baseline,sec:locality})
  \item What do alignment, loop tails, stride and gather access, and data type cost? (\cref{sec:alignment,sec:stride,sec:dtype})
  \item Where do these kernels sit in a roofline model of this machine? (\cref{sec:roofline})
  \item When does the expected SIMD benefit fail to appear? (\cref{sec:negative})
\end{enumerate}

In brief: SIMD helps most when data fits in L1/L2, where \code{native} gives roughly $5.8\times$ for
SAXPY and the cache feeds the vector units fast enough to approach the lane count; the advantage
shrinks step-by-step as the working set crosses L2 and L3 and falls to $1.88\text{--}1.99\times$ for SAXPY
in DRAM ($1.5\text{--}1.6\times$ multiply, $2.6\text{--}2.8\times$ stencil), where the vector build runs at the memory-bandwidth roof and only the scalar build's own
instruction cost keeps the ratio above $1\times$. Data type matters through lane count (float64
roughly halves the gain) and access pattern matters even more (any stride~$\ge2$ removes the SIMD
advantage; gather, even with AVX2 \code{vgatherdps}, never gains). Two cases show the limits of ``it vectorized, so it
must be faster'': memory-bound SAXPY (\cref{sec:negative}) and the dot-product reduction
(\cref{sec:vectorization}).

% !TeX root = ../main.tex
\section{Methodology}\label{sec:method}
\rubric{Baseline, correctness, reproducibility}{15}{Reference validation, controlled timing,
reproducible scripts and a documented environment (with \cref{sec:correctness,sec:baseline}).}

\subsection{Platform}
\newcommand{\acspinning}{Each run is pinned to one physical P-core
(\code{--cpu 2}, \code{sched_setaffinity}); its SMT sibling is left idle and no other
pinning/isolation is used.}
%% Measurement environment table (common requirement: measurement discipline).
%% Values come from common/system_info.tex. Before \input-ing this file a project may define
%%   \acspinning   how runs were pinned/isolated
%%   \acsenvextra  extra rows, e.g.  \newcommand{\acsenvextra}{SSD & model, firmware, PCIe link \\}
\providecommand{\acspinning}{\todo{e.g.\ \code{taskset -c 2} on a P-core; isolcpus if used}}
\providecommand{\acsenvextra}{}
\begin{table}[tbp]
  \centering\small
  \caption{Measurement environment.}\label{tab:environment}
  \begin{tabularx}{\linewidth}{@{}>{\raggedright\arraybackslash}p{0.2\linewidth}X@{}}
    \toprule
    CPU               & \sysCPU{} (\sysArch) \\
    SIMD ISA          & \sysISA \\
    Topology          & \sysCores; NUMA: \sysNUMA \\
    SMT               & \sysSMT \\
    Caches            & \sysCaches \\
    Memory            & \sysMemory; DIMMs: \sysDIMM \\
    Pages             & base \sysPageSize; THP: \sysHugePages \\
    OS / kernel       & \sysOS; \sysKernel \\
    Frequency policy  & governor: \sysGovernor; driver: \sysDriver; turbo: \sysTurbo; EPP: \sysEPP \\
    Pinning / isolation & \acspinning; \sysIsolation \\
    Compilers         & \sysGCC; \sysClang \\
    Tools             & \sysPerf{} (perf\_event\_paranoid \sysPerfParanoid); \sysFio; MLC: \sysMLC \\
    GPU (optional)    & \sysGPU; CUDA \sysCUDA \\
    \acsenvextra
    Collected         & \sysCollected \\
    \bottomrule
  \end{tabularx}
\end{table}


This CPU (Comet Lake, client Skylake core) implements SSE/AVX/AVX2 and FMA but \emph{not}
AVX-512, so the widest usable vector register is \qty{256}{bit} (\code{ymm}): \num{8} lanes of
float32 or \num{4} of float64. With \code{-march=native} GCC targets AVX2+FMA and prefers
\qty{256}{bit} vectors; the default (generic \code{x86-64}) target only guarantees SSE2, so the
plain auto-vectorized build emits \qty{128}{bit} (\code{xmm}) code. That 256-vs-128-bit gap is the
main reason we keep \code{autovec} and \code{native} as separate variants (\cref{tab:variants}).

\subsection{Kernels}
\begin{table}[tbp]
  \centering\small
  \caption{Kernels and floating-point operations (FLOPs) per processed element.}\label{tab:kernels}
  \begin{tabular}{@{}lll@{}}
    \toprule
    Kernel & Definition & FLOPs / element \\
    \midrule
    SAXPY               & $y_i \leftarrow a\,x_i + y_i$                         & 2 \\
    Dot product         & $s \leftarrow \sum_i x_i\, y_i$                        & 2 \\
    Elementwise multiply & $z_i \leftarrow x_i\, y_i$                            & 1 \\
    3-point stencil     & $y_i \leftarrow a\,x_{i-1} + b\,x_i + c\,x_{i+1}$      & 5 \\
    \bottomrule
  \end{tabular}
\end{table}
We count a multiply and an add as one FLOP each, so a fused multiply--add is two FLOPs; SAXPY and
the dot product are therefore 2~FLOP/element (one FMA), elementwise multiply is 1, and the stencil
is 5 (three multiplies, two adds). All four kernels from the project brief are implemented and
validated; SAXPY is used as the representative kernel for the locality, alignment and stride
studies because it has the simplest unit-stride access and a well-known arithmetic intensity.

\subsection{Build variants}
\begin{table}[tbp]
  \centering\small
  \caption{Build variants. Only the kernel translation unit (\code{src/kernels.c}) differs; the
  timing harness (\code{main.c}, \code{common.c}) and the reference (\code{reference.c}) are
  compiled identically across variants.}\label{tab:variants}
  \begin{tabularx}{\linewidth}{@{}lX@{}}
    \toprule
    Variant & Kernel compiler flags (from \code{simd_bench --info}) \\
    \midrule
    scalar    & \code{-O3 -fno-tree-vectorize -fno-tree-slp-vectorize} \\
    autovec   & \code{-O3} (generic \code{x86-64} $\Rightarrow$ SSE2, \qty{128}{bit}) \\
    native    & \code{-O3 -march=native} (AVX2+FMA, \qty{256}{bit}) \\
    native512 & not available on this CPU (no AVX-512; \qty{256}{bit} is the widest ISA) \\
    \bottomrule
  \end{tabularx}
\end{table}
The primary compiler is GCC \sysGCC{}. The reference implementation additionally uses
\code{-O2 -fno-tree-vectorize -ffp-contract=off} so it evaluates each expression in strict,
non-fused order and stays a trustworthy oracle. We deliberately do \emph{not} pass
\code{-ffast-math} to any variant: this keeps float results reproducible but, as
\cref{sec:vectorization,sec:negative} show, it also prevents the compiler from reassociating the
dot-product reduction---an effect we wanted to measure rather than hide.

\subsection{Timing and statistics}
\begin{itemize}
  \item Timer: \code{clock_gettime(CLOCK_MONOTONIC_RAW)} (resolution \qty{1}{\nano\second} on this
        host, from \code{simd_bench --info}). Each repetition times a calibrated batch of kernel
        invocations lasting at least \qty{25}{\milli\second}, i.e.\ roughly $2.5\times10^{7}$ timer
        ticks, so timer resolution and call overhead are negligible.
  \item Warm-up: 3 untimed invocations per configuration; repetitions: 12 per configuration.
  \item Reported statistic: median; error bars show the interquartile range (IQR) across the 12
        repetitions. The median resists the occasional scheduler/turbo outlier better than the mean.
  \item Pinning, SMT, governor and turbo: \cref{tab:environment}. Turbo is left enabled and the
        governor is \sysGovernor{}; the resulting frequency variation is a documented limitation
        (\cref{sec:anomalies}).
  \item Every kernel result is written through a \code{volatile} sink (\code{acs_consume_*}) and the
        kernels live in their own translation unit, so no loop is eliminated as dead code.
  \item Configuration order: the sweep is run as nested loops over size and variant
        (\code{scripts/run.sh}); it is not shuffled, which we note as a limitation.
\end{itemize}

\subsection{Data initialization and working-set sizes}
Inputs are filled with a deterministic linear-congruential sequence in $[-0.5, 0.5)$ (fixed seed),
which avoids denormals, NaNs and overflow; SAXPY's in-place update grows $y$ only linearly in the
repetition count, so it never overflows float32. The gather pattern uses a Fisher--Yates
permutation of $[0,N)$ so every element is touched exactly once with no locality.

Each kernel streams two or three arrays, so the bytes touched are (arrays)$\times$(element
size)$\times N$. For float32 SAXPY (two arrays, \qty{8}{\byte}/element) the working set crosses
this machine's \qty{32}{\kibi\byte} L1d near $N\approx\num{4096}$, the \qty{256}{\kibi\byte} L2 near
$N\approx\num{32768}$, and the \qty{20}{\mebi\byte} L3 near $N\approx\num{2.6e6}$; beyond that it is
DRAM-resident. These boundaries are annotated on the locality plot (\cref{sec:locality}).

% !TeX root = ../main.tex
\section{Correctness Validation}\label{sec:correctness}

Every build variant is checked two independent ways before any of its timings are trusted
(\code{make test}, \code{scripts/validate.sh}).

\emph{Reference comparison.} \code{src/reference.c} contains plain scalar loops compiled with
\code{-fno-tree-vectorize -ffp-contract=off}, so it evaluates each kernel in strict, non-fused,
left-to-right order. The driver runs the kernel under test and the reference on identical inputs and
compares element-by-element (or the scalar result, for the dot product) before timing.

\begin{knobs}
  Cases  & kernels $\times$ \{float32, float64\}; $N$ around multiples of the lane count
           ($1,2,3,7,8,9,15,16,17,31,33,64,127,1000,1024,4095,4096,4097,65536$); start offsets
           $4\ldots\qty{60}{\byte}$; strides $2,4,8,16$; gather; int32 multiply \\
  Tolerance & relative error $\le 10^{-5}$ (float32) and $\le 10^{-12}$ (float64). The bound is loose
             enough for the one legal source of disagreement---FMA contraction in the vectorized
             builds versus the non-fused reference---and tight enough to catch real bugs. \\
  Command & \code{make test} \\
\end{knobs}

\emph{Known-answer self-test} (\code{simd_bench --selftest}), independent of the reference: the dot
product of two all-ones vectors equals $N$; $y=2x+y$ with $x=1,y=3$ gives $5$ everywhere; $z=x\,y$
with $x=2,y=3$ gives $6$; and the stencil with $a=b=c=1$ on a constant-1 input gives $3$ in the
interior. These guard against a reference and kernel sharing the same mistake.

The sweep covers sizes with and without a vector-width remainder (to exercise the scalar/masked tail)
and deliberately misaligned start addresses (to exercise the loop prologue), so correctness is
established precisely in the regimes the later experiments stress.

Result: all \num{570} validation cases pass for every variant (scalar, autovec, native); the maximum
observed float32 relative error was below $10^{-6}$, consistent with FMA contraction alone.

% !TeX root = ../main.tex
\section{Vectorization Verification}\label{sec:vectorization}
\rubric{Vectorization verification}{15}{Clear evidence that SIMD instructions are used, with a correct
reading of vector width, FMA and masking.}

\begin{hypothesis}
  \todo{Which loops you expect each variant to vectorize, at what width, and which you expect the
  compiler to reject (and why).}
\end{hypothesis}

\Cref{tab:vec-summary} counts, for each kernel function, the instructions that operate on vector
registers or use packed operations (\code{make report}).

\begin{table}[tbp]
  \centering\small
  \caption{SIMD instructions per kernel function and build variant.}\label{tab:vec-summary}
  \IfDataExists{vectorization}{%
    \pgfplotstabletypeset[
      col sep=comma,
      columns={variant,function,insns,simd_insns,widest_register},
      columns/variant/.style={column name=Variant, string type},
      columns/function/.style={column name=Function, string type, string replace*={_}{\_}},
      columns/insns/.style={column name=Instructions},
      columns/simd_insns/.style={column name=SIMD instructions},
      columns/widest_register/.style={column name=Widest register, string type},
      every head row/.style={before row=\toprule, after row=\midrule},
      every last row/.style={after row=\bottomrule},
    ]{data/vectorization.csv}%
  }
\end{table}

\subsection{Compiler report}
\begin{lstlisting}[style=excerpt]
TODO: 2-5 relevant lines from data/raw/vectorization/<run>/report_<variant>.txt
\end{lstlisting}

\subsection{Disassembly}
\begin{lstlisting}[style=excerpt]
TODO: 5-10 lines of a vectorized loop body, e.g. disasm/native/saxpy_f32.s
\end{lstlisting}

\todo{Interpret the evidence: register width (xmm/ymm/zmm) and lanes per element type; FMA
(\code{vfmadd*}); masked tails (k-mask registers) versus a scalar epilogue; loops that did not
vectorize and the compiler's reason.}

% !TeX root = ../main.tex
\section{Scalar vs.\ SIMD Baseline}\label{sec:baseline}
\rubric{Baseline, correctness, reproducibility}{15}{Speedup and achieved throughput for every kernel
across working-set sizes.}

\begin{hypothesis}
  \todo{Expected speedup per kernel relative to the lane count, and how it should change between
  cache-resident and DRAM-resident $N$.}
\end{hypothesis}

\begin{knobs}
  Fixed       & float32; unit stride; aligned arrays; one thread pinned to CPU \tbd \\
  Varied      & kernel; variant (scalar, autovec, native); $N = 2^{8} \ldots 2^{26}$ \\
  Repetitions & \tbd{} per point; median, IQR error bars \\
  Script      & \code{SIMD_Advantage_Profiling/scripts/run.sh baseline} \\
\end{knobs}

\begin{figure}[tbp]
  \centering
  \IfDataExists{baseline}{%
    \begin{tikzpicture}
      \begin{axis}[acs, xmode=log, log basis x=2, xlabel={$N$ (elements)}, ylabel={Speedup over scalar},
                   legend pos=outer north east, width=0.75\linewidth]
        \addplot table[x=n, y=speedup_vs_scalar, only if={kernel}{saxpy}, only if={variant}{autovec}] {data/baseline.csv};
        \addlegendentry{SAXPY autovec}
        \addplot table[x=n, y=speedup_vs_scalar, only if={kernel}{saxpy}, only if={variant}{native}] {data/baseline.csv};
        \addlegendentry{SAXPY native}
        % TODO: remaining kernels; cache boundaries, e.g. \acsvline{<L1d elements>}{L1d}
      \end{axis}
    \end{tikzpicture}%
  }
  \caption{Speedup of the SIMD builds over the scalar build. \todo{one-sentence takeaway}}\label{fig:baseline-speedup}
\end{figure}

\begin{figure}[tbp]
  \centering
  \IfDataExists{baseline}{%
    \begin{tikzpicture}
      \begin{axis}[acs, xmode=log, log basis x=2, ymode=log, xlabel={$N$ (elements)}, ylabel={Throughput (elements/s)},
                   legend pos=outer north east, width=0.75\linewidth]
        \addplot+[iqr bars] table[x=n, y=elems_per_s_median, iqr of=elems_per_s, only if={kernel}{saxpy}, only if={variant}{scalar}] {data/baseline.csv};
        \addlegendentry{SAXPY scalar}
        \addplot+[iqr bars] table[x=n, y=elems_per_s_median, iqr of=elems_per_s, only if={kernel}{saxpy}, only if={variant}{native}] {data/baseline.csv};
        \addlegendentry{SAXPY native}
        % TODO: report GFLOP/s instead: y expr={\thisrow{elems_per_s_median} * <FLOPs per element> / 1e9}
      \end{axis}
    \end{tikzpicture}%
  }
  \caption{Achieved throughput. \todo{one-sentence takeaway}}\label{fig:baseline-throughput}
\end{figure}

\subsection{Analysis}
\todo{Speedups versus lane count per kernel; why reductions (dot) and stencils behave differently from
SAXPY and elementwise multiply; where the curves bend and why.}

% !TeX root = ../main.tex
\section{Locality Sweep}\label{sec:locality}
\rubric{Locality and memory-bound analysis}{20}{Working-set transitions identified, and the compression
of SIMD gains in the memory-bound regime explained.}

\begin{hypothesis}
  \todo{At which $N$ each cache transition should appear for this kernel, and how scalar and SIMD
  builds should diverge or converge beyond LLC.}
\end{hypothesis}

\begin{knobs}
  Fixed       & kernel \tbd; float32; unit stride; aligned; pinned to CPU \tbd \\
  Varied      & $N$ (dense grid across L1/L2/LLC/DRAM); variant \\
  Counters    & \todo{perf events recorded (\code{PERF=1}), e.g.\ LLC-load-misses} \\
  Script      & \code{SIMD_Advantage_Profiling/scripts/run.sh locality} \\
\end{knobs}

\begin{figure}[tbp]
  \centering
  \IfDataExists{locality}{%
    \begin{tikzpicture}
      \begin{axis}[acs, xmode=log, log basis x=2, ymode=log, xlabel={$N$ (elements)}, ylabel={Time per element (ns)},
                   legend pos=north west]
        \addplot+[iqr bars] table[x=n, y=ns_per_elem_median, iqr of=ns_per_elem, only if={variant}{scalar}] {data/locality.csv};
        \addlegendentry{scalar}
        \addplot+[iqr bars] table[x=n, y=ns_per_elem_median, iqr of=ns_per_elem, only if={variant}{autovec}] {data/locality.csv};
        \addlegendentry{autovec}
        \addplot+[iqr bars] table[x=n, y=ns_per_elem_median, iqr of=ns_per_elem, only if={variant}{native}] {data/locality.csv};
        \addlegendentry{native}
        % TODO: mark cache boundaries in elements: \acsvline{<N>}{L1d} \acsvline{<N>}{L2} \acsvline{<N>}{LLC}
      \end{axis}
    \end{tikzpicture}%
  }
  \caption{Time per element across cache levels. \todo{takeaway}}\label{fig:locality}
\end{figure}

\subsection{Analysis}
\todo{Where each transition occurs (bytes = arrays $\times$ element size $\times N$) and what the
counters show there; why the SIMD speedup compresses once memory traffic dominates.}

% !TeX root = ../main.tex
\section{Alignment and Tail Handling}\label{sec:alignment}
\rubric{Alignment, tail, stride/gather studies}{20}{Quantified alignment, prologue/epilogue and masking
effects with a convincing explanation (with \cref{sec:stride}).}

\begin{hypothesis}
  A start address that is not \qty{32}{\byte}-aligned should force the \qty{256}{bit} loads to be
  unaligned and/or trigger a scalar peel prologue, costing a modest fixed overhead that is most
  visible at small $N$. Because this is a client Skylake core, unaligned \code{vmovups} that happen
  to stay within a cache line should be nearly free, so we expect the penalty to be small
  (tens of percent) and to recur with \qty{64}{\byte} period as the offset sweeps a cache line.
\end{hypothesis}
\begin{hypothesis}
  (Added for the extended sweep, before it was run.) The \code{scalar} build issues \qty{4}{\byte}
  loads that never straddle a line, so it should be insensitive to the offset; \code{autovec}
  (\qty{16}{\byte} loads) should show a smaller penalty than \code{native}. With a DRAM-resident
  array the misalignment cost should disappear into the memory stall time for all builds. For short
  loops, the time per \emph{call} should grow in steps of \num{8} elements for \code{native} and
  \num{4} for \code{autovec} (one vector iteration per step), with the vectorized epilogue making the
  remainder cheap but not free.
\end{hypothesis}

\begin{knobs}
  Fixed       & kernel SAXPY; float32; unit stride \\
  Varied      & start offset $0\ldots\qty{60}{\byte}$ (multiples of \num{4}) at $N=\num{4096}$
                (L1-resident, isolates the prologue from bandwidth) and $N=\num{16777216}$
                (\qty{128}{\mebi\byte}, DRAM-resident); tail: every $N$ from \num{8} to \num{72} and
                $N$ stepping across an \num{8}-lane boundary near \num{4096}, offset $0$; all three
                build variants \\
  Script      & \code{SIMD_Advantage_Profiling/scripts/run.sh alignment} \\
\end{knobs}

% New layout (all variants, DRAM offsets, short-N tail) once the extended sweep has produced
% data/alignment_tail.csv; until then the original native-only figure from alignment.csv.
\IfFileExists{data/alignment_tail.csv}{%
\begin{figure}[tbp]
  \centering
  \begin{subfigure}[t]{0.49\linewidth}
    \begin{tikzpicture}
      \begin{axis}[acs, height=0.85\linewidth, xlabel={Start offset (bytes)}, ylabel={Time per element (ns)},
                   legend style={font=\scriptsize}, legend pos=north east, ymin=0]
        \acsvariantlegend
        \addplot+[var scalar, forget plot, iqr bars] table[x=offset, y=ns_per_elem_median, iqr of=ns_per_elem, only if={variant}{scalar}, only if={n}{4096}] {data/alignment.csv};
        \addplot+[var autovec, forget plot, iqr bars] table[x=offset, y=ns_per_elem_median, iqr of=ns_per_elem, only if={variant}{autovec}, only if={n}{4096}] {data/alignment.csv};
        \addplot+[var native, forget plot, iqr bars] table[x=offset, y=ns_per_elem_median, iqr of=ns_per_elem, only if={variant}{native}, only if={n}{4096}] {data/alignment.csv};
      \end{axis}
    \end{tikzpicture}
    \caption{$N=\num{4096}$ (L1-resident).}
  \end{subfigure}\hfill
  \begin{subfigure}[t]{0.49\linewidth}
    \begin{tikzpicture}
      \begin{axis}[acs, height=0.85\linewidth, xlabel={Start offset (bytes)}, ylabel={Time per element (ns)}, ymin=0]
        \addplot+[var scalar, forget plot, iqr bars] table[x=offset, y=ns_per_elem_median, iqr of=ns_per_elem, only if={variant}{scalar}, only if={n}{16777216}] {data/alignment.csv};
        \addplot+[var autovec, forget plot, iqr bars] table[x=offset, y=ns_per_elem_median, iqr of=ns_per_elem, only if={variant}{autovec}, only if={n}{16777216}] {data/alignment.csv};
        \addplot+[var native, forget plot, iqr bars] table[x=offset, y=ns_per_elem_median, iqr of=ns_per_elem, only if={variant}{native}, only if={n}{16777216}] {data/alignment.csv};
      \end{axis}
    \end{tikzpicture}
    \caption{$N=\num{16777216}$ (DRAM-resident).}
  \end{subfigure}
  \caption{Cost of a misaligned start for SAXPY, by build variant (median, IQR error bars).}\label{fig:alignment}
\end{figure}

\begin{figure}[tbp]
  \centering
  \begin{subfigure}[t]{0.49\linewidth}
    \begin{tikzpicture}
      \begin{axis}[acs, height=0.85\linewidth, xlabel={$N$ (elements)}, ylabel={Time per call (ns)},
                   restrict x to domain=0:100, xmin=6, xmax=74, xtick={8,16,...,72}, legend style={font=\scriptsize}, legend pos=north west]
        \acsvariantlegend
        \addplot+[mark size=1pt, var scalar, forget plot] table[x=n, y expr=\thisrow{ns_per_elem_median}*\thisrow{n}, only if={variant}{scalar}] {data/alignment_tail.csv};
        \addplot+[mark size=1pt, var autovec, forget plot] table[x=n, y expr=\thisrow{ns_per_elem_median}*\thisrow{n}, only if={variant}{autovec}] {data/alignment_tail.csv};
        \addplot+[mark size=1pt, var native, forget plot] table[x=n, y expr=\thisrow{ns_per_elem_median}*\thisrow{n}, only if={variant}{native}] {data/alignment_tail.csv};
      \end{axis}
    \end{tikzpicture}
    \caption{Short loops: time per call, $N=8\ldots72$.}
  \end{subfigure}\hfill
  \begin{subfigure}[t]{0.49\linewidth}
    \begin{tikzpicture}
      \begin{axis}[acs, height=0.85\linewidth, xlabel={$N$ (elements)}, ylabel={Time per element (ns)},
                   restrict x to domain=4080:4140, xmin=4084, xmax=4132, ymin=0]
        \addplot+[var scalar, forget plot, iqr bars] table[x=n, y=ns_per_elem_median, iqr of=ns_per_elem, only if={variant}{scalar}] {data/alignment_tail.csv};
        \addplot+[var autovec, forget plot, iqr bars] table[x=n, y=ns_per_elem_median, iqr of=ns_per_elem, only if={variant}{autovec}] {data/alignment_tail.csv};
        \addplot+[var native, forget plot, iqr bars] table[x=n, y=ns_per_elem_median, iqr of=ns_per_elem, only if={variant}{native}] {data/alignment_tail.csv};
      \end{axis}
    \end{tikzpicture}
    \caption{Across an \num{8}-lane boundary near $N=\num{4096}$.}
  \end{subfigure}
  \caption{Tail handling: remainder elements processed by the vectorized epilogue (offset $0$).}\label{fig:tail}
\end{figure}
}{%
\begin{figure}[tbp]
  \centering
  \IfDataExists{alignment}{%
    \begin{subfigure}[t]{0.49\linewidth}
      \begin{tikzpicture}
        \begin{axis}[acs, height=0.8\linewidth, xlabel={Start offset (bytes)}, ylabel={Time per element (ns)}]
          \addplot+[var native, forget plot, iqr bars] table[x=offset, y=ns_per_elem_median, iqr of=ns_per_elem, only if={variant}{native}, only if={n}{4096}] {data/alignment.csv};
        \end{axis}
      \end{tikzpicture}
      \caption{Misaligned starts ($N=4096$).}
    \end{subfigure}\hfill
    \begin{subfigure}[t]{0.49\linewidth}
      \begin{tikzpicture}
        \begin{axis}[acs, height=0.8\linewidth, xlabel={$N$ (elements)}, ylabel={Time per element (ns)}]
          \addplot+[var native, forget plot, iqr bars] table[x=n, y=ns_per_elem_median, iqr of=ns_per_elem, only if={variant}{native}, only if={offset}{0}] {data/alignment.csv};
        \end{axis}
      \end{tikzpicture}
      \caption{Tail length (offset $0$).}
    \end{subfigure}%
  }
  \caption{Alignment and tail-handling costs for native SAXPY.}\label{fig:alignment}\label{fig:tail}
\end{figure}
}


\subsection{Analysis}
\emph{Misaligned starts, L1-resident} (\cref{fig:alignment}a). The penalty is real, it scales with
vector width, and it repeats with the load width. Native SAXPY runs at \qty{0.099}{\nano\second} per
element at offsets $0$ and \qty{32}{\byte} (both \qty{32}{\byte}-aligned, since the base is
\qty{64}{\byte}-aligned). Offsets \numrange{4}{24} cost \qty{0.143}{\nano\second} (\SI{+45}{\percent}),
and \numlist{28;36;60} cost \qty{0.157}{\nano\second} (\SI{+59}{\percent}). \code{autovec}
(\qty{16}{\byte} loads) is fast at offsets $0$, \num{16} and \num{32} (\qtyrange{0.190}{0.196}{\nano\second}),
i.e.\ with the \qty{16}{\byte} period of its loads, and up to \SI{26}{\percent} slower elsewhere.
\code{scalar} varies by only \SI{1.8}{\percent} (\qtyrange{0.563}{0.573}{\nano\second}). This
ordering, native \SI{60}{\percent} > autovec \SI{26}{\percent} > scalar \SI{2}{\percent}, is what
cache-line splits predict. At a non-multiple-of-32 offset every second \qty{32}{\byte} load straddles
a \qty{64}{\byte} line, but only every fourth \qty{16}{\byte} load does, and a \qty{4}{\byte} scalar
load never does. GCC does not peel the loop to an aligned start, so every split load pays.

\emph{DRAM-resident} (\cref{fig:alignment}b). At $N=\num{16777216}$ the native spread across
offsets shrinks to \SI{2.0}{\percent} (\qtyrange{0.411}{0.419}{\nano\second}) and autovec's to
\SI{5.5}{\percent}: the extra split-load cycles hide under the DRAM misses that dominate each
iteration, so for the vector builds alignment is a cache-resident concern only. Scalar shows a
\SI{18}{\percent} spread, but it is bimodal (\num{0.713} or \qty{0.84}{\nano\second}) and does
not follow the offset: offsets \num{8}, \num{24} and \num{32} are fast while \num{4}, \num{16} and
\num{20} are slow, although a \qty{4}{\byte} load is never split. Each offset is a separate
process with freshly allocated \qty{128}{\mebi\byte} arrays, so we read this as process-to-process
variation (see the wide scalar IQR at the same $N$ in \cref{fig:baseline-throughput}), not as an
alignment effect.

\emph{Tails} (\cref{fig:tail}). Near $N=\num{4096}$ the time per element is flat across the
\num{8}-lane boundary (native \qtyrange{0.098}{0.105}{\nano\second}, autovec
\qtyrange{0.184}{0.198}{\nano\second}, scalar \qtyrange{0.57}{0.63}{\nano\second}). With
\numrange{1}{7} leftover elements out of \num{4096}, the remainder is invisible. Short loops expose
it. For $N=8\ldots72$, native's time per call is lowest at multiples of \num{8} (\qty{13.6}{\nano\second}
at $N=16$ rising to \qty{17.4}{\nano\second} at $N=72$, only $\approx\qty{0.07}{\nano\second}$ per
element) and highest at $N\equiv7 \pmod 8$, \qtyrange{2}{3}{\nano\second} more. That sawtooth is
the remainder path from the disassembly (\cref{sec:vectorization}): a \qty{128}{bit} epilogue for
\num{4} elements and then up to three scalar iterations. \code{autovec} grows at
$\approx\qty{0.15}{\nano\second}$ per element and \code{scalar} linearly at
$\approx\qty{0.63}{\nano\second}$ per element, both with occasional single-point outliers. Every
call also pays a fixed $\approx\qtyrange{10}{13}{\nano\second}$ (call, loop setup and the run-time
aliasing check of the versioned loop) before the first element. So the per-element speedup of
short loops is far below the $5.8\times$ at $N=\num{4096}$: native beats scalar by only
$1.11\times$ at $N=8$ and by $3.20\times$ at $N=72$. There are no AVX-512 mask registers here
(no AVX-512 on this CPU); a masked tail would be the AVX-512 alternative to this narrowing epilogue.

% !TeX root = ../main.tex
\section{Stride and Gather Access}\label{sec:stride}

\begin{hypothesis}
  Unit stride lets each \qty{64}{\byte} cache line serve \num{16} consecutive float32 elements, so
  SIMD helps. As the stride grows, each vector still needs \num{8} elements but they come from
  further apart: once the stride reaches \num{16} elements (\qty{64}{\byte}) every element sits in its
  own cache line and effective bandwidth drops by up to $16\times$. We therefore expect time per
  processed element to rise roughly linearly with stride until it saturates at one cache line per
  element, and we expect the SIMD advantage to vanish for any non-unit stride, because the compiler
  must gather lanes individually. Gather (random permutation) should be the worst case.
\end{hypothesis}

\begin{knobs}
  Fixed       & SAXPY; float32; native and scalar builds \\
  Varied      & access (unit, stride $1/2/4/8/16$, gather); $N$ cache-resident (\num{32768}) and
                DRAM-resident (\num{4194304}) \\
  Definition  & stride $s$ processes elements $0, s, 2s, \ldots$ over an array of $N s$ elements;
                gather applies SAXPY through a random permutation index so every element is touched
                once with no spatial locality. \\
  Script      & \code{SIMD_Advantage_Profiling/scripts/run.sh stride} \\
\end{knobs}

\begin{figure}[tbp]
  \centering
  \IfDataExists{stride}{%
    \begin{tikzpicture}
      \begin{axis}[acs, xmode=log, log basis x=2, xlabel={Stride (elements)}, ylabel={Time per processed element (ns)}, legend pos=north west]
        \addplot+[iqr bars] table[x=stride, y=ns_per_elem_median, iqr of=ns_per_elem, only if={access}{stride}, only if={variant}{scalar}, only if={n}{4194304}] {data/stride.csv};
        \addlegendentry{scalar (DRAM)}
        \addplot+[iqr bars] table[x=stride, y=ns_per_elem_median, iqr of=ns_per_elem, only if={access}{stride}, only if={variant}{native}, only if={n}{4194304}] {data/stride.csv};
        \addlegendentry{native (DRAM)}
      \end{axis}
    \end{tikzpicture}%
  }
  \caption{Effect of stride on native and scalar SAXPY time per \emph{processed} element, DRAM-resident
  ($N=\num{4194304}$). The two curves coincide for every non-unit stride.}\label{fig:stride}
\end{figure}

\begin{table}[tbp]
  \centering\small
  \caption{Time per processed element (ns), DRAM-resident SAXPY, $N=\num{4194304}$.}\label{tab:stride}
  \begin{tabular}{@{}lrrr@{}}
    \toprule
    Access & scalar & native & native speedup \\
    \midrule
    unit (stride 1)   & 0.453 & 0.252 & $1.80\times$ \\
    stride 2          & 0.790 & 0.787 & $1.00\times$ \\
    stride 4          & 1.540 & 1.509 & $1.02\times$ \\
    stride 8          & 2.944 & 2.977 & $0.99\times$ \\
    stride 16         & 6.074 & 6.122 & $0.99\times$ \\
    gather (random)   & 10.83 & 10.69 & $1.01\times$ \\
    \bottomrule
  \end{tabular}
\end{table}

\subsection{Analysis}
\Cref{tab:stride} makes the pattern stark. Only unit stride lets \code{native} beat \code{scalar}
(\,$1.80\times$\,); at stride~2 and beyond the two builds are within measurement noise of each other.
Two things happen at once. First, the effective bandwidth collapses: time per processed element
roughly doubles from stride~2 to~4 to~8 to~16 (\num{0.79} $\to$ \num{1.51} $\to$ \num{2.98}
$\to$ \qty{6.12}{\nano\second}), because once the stride reaches \num{16} elements each access pulls a
fresh \qty{64}{\byte} line and $15/16$ of every line is wasted. Second, the SIMD advantage
disappears: a strided or gathered load cannot be a single contiguous \code{vmovups}, so the compiler
either falls back to scalar element loads or (for gather) emits per-lane address arithmetic; the
vector ALU is then starved regardless of width. The disassembly confirms this---\code{native}'s
gather kernel stays on \qty{128}{bit} \code{xmm} and never reaches \code{ymm}
(\cref{tab:vec-summary}), and GCC does not emit an AVX2 \code{vgather} here at all. Gather is the
worst case at \qty{10.7}{\nano\second} per element, $\approx\!42\times$ slower than unit stride,
because the random permutation defeats the hardware prefetcher as well. The practical lesson is that
access pattern, not vector width, is the first-order knob for these kernels.

% !TeX root = ../main.tex
\section{Data Type Comparison}\label{sec:dtype}

\begin{hypothesis}
  \todo{Expected float32 vs.\ float64 (and int32) speedups from lane count, operation cost and bytes
  per element, in cache and in DRAM.}
\end{hypothesis}

\begin{knobs}
  Fixed       & unit stride; aligned \\
  Varied      & data type; kernel; variant; $N$ \\
  Script      & \code{SIMD_Advantage_Profiling/scripts/run.sh dtype} \\
\end{knobs}

\begin{figure}[tbp]
  \centering
  \IfDataExists{dtype}{%
    \begin{tikzpicture}
      \begin{axis}[acs, xmode=log, log basis x=2, xlabel={$N$ (elements)}, ylabel={Speedup over scalar}, legend pos=north east]
        \addplot table[x=n, y=speedup_vs_scalar, only if={kernel}{mul}, only if={variant}{native}, only if={dtype}{f32}] {data/dtype.csv};
        \addlegendentry{multiply float32}
        \addplot table[x=n, y=speedup_vs_scalar, only if={kernel}{mul}, only if={variant}{native}, only if={dtype}{f64}] {data/dtype.csv};
        \addlegendentry{multiply float64}
        % TODO: other kernels; optional int32
      \end{axis}
    \end{tikzpicture}%
  }
  \caption{Speedup by data type. \todo{takeaway}}\label{fig:dtype}
\end{figure}

\subsection{Analysis}
\todo{Explain with lane count, arithmetic intensity (FLOPs per byte) and per-operation cost; note where
float64 halves the lanes but not necessarily the speedup.}

% !TeX root = ../main.tex
\section{Roofline Interpretation}\label{sec:roofline}
\rubric{Roofline and quantitative insight}{20}{Correct arithmetic intensity, use of measured memory
bandwidth, and a justified compute- vs.\ memory-bound conclusion.}

\begin{hypothesis}
  All four kernels have very low arithmetic intensity (well under \qty{1}{\flop\per\byte}), so on a
  machine whose single-core bandwidth is tens of \unit{\giga\byte\per\second} and whose peak vector
  rate is over \qty{100}{\giga\flop\per\second} they should all sit on the \emph{bandwidth} slope of
  the roofline when DRAM-resident, i.e.\ be memory-bound. The stencil, with the highest intensity,
  should come closest to the compute roof.
\end{hypothesis}
\begin{hypothesis}
  (For the measured ceilings.) Even the L1 bandwidth roof should put the
  ridge point above \qty{0.5}{\flop\per\byte}, so the cache-resident native kernels should also be
  bounded by bandwidth (or by the load/store ports that define it), not by the FMA units.
\end{hypothesis}

Attainable performance at arithmetic intensity $I$ (FLOPs per byte of memory traffic) is bounded by
\begin{equation}
  P(I) = \min\bigl(P_{\text{peak}},\; B_{\ell} \cdot I\bigr),
\end{equation}
where $P_{\text{peak}}$ is the core's peak FLOP rate and $B_{\ell}$ the bandwidth of the memory level
$\ell$ the data lives in \cite{williams2009roofline}. We draw one bandwidth roof per level (a
cache-aware roofline). Both ceilings are \emph{measured by separate microbenchmarks}
(\cref{sec:ceilings-method}), never derived from the kernels plotted against them.

\begin{table}[tbp]
  \centering\small
  \caption{Measured roofline ceilings, single core (\code{run.sh ceilings}). Top: register-only
  accumulator loops (\num{12} independent chains); FLOP/cycle and clock from the user-space cycle
  counter. Bottom: STREAM-style bandwidth at one footprint per level (GiB/s). \enquote{incl.\ WA}
  adds the write-allocate read; the roofs in \cref{fig:roofline} use the explicit triad in L1
  (stores hit L1, no write-allocate) and triad incl.\ WA below it. The read stream ($^{*}$) is
  add-latency-bound (\cref{sec:ceilings-method}) and is not used as a roof.}\label{tab:ceilings}
  \IfDataExists{ceilings_peak}{%
    \pgfplotstabletypeset[
      col sep=comma,
      columns={isa,gflops_median,flops_per_cycle,ghz},
      columns/isa/.style={column name=Peak loop, string type, string replace*={_}{\_}},
      columns/gflops_median/.style={column name=GFLOP/s, fixed, zerofill, precision=1},
      columns/flops_per_cycle/.style={column name=FLOP/cycle, fixed, zerofill, precision=2},
      columns/ghz/.style={column name=GHz, fixed, zerofill, precision=2},
      every head row/.style={before row=\toprule, after row=\midrule},
      every last row/.style={after row=\bottomrule},
    ]{data/ceilings_peak.csv}%
  }\par\medskip
  \IfDataExists{ceilings_bw}{%
    \pgfplotstabletypeset[
      col sep=comma,
      columns={level,bytes,read,copy,triad,triad_total},
      columns/level/.style={column name=Level, string type},
      columns/bytes/.style={column name=Footprint (B), int detect},
      columns/read/.style={column name=read$^{*}$, fixed, zerofill, precision=1},
      columns/copy/.style={column name=copy, fixed, zerofill, precision=1},
      columns/triad/.style={column name=triad, fixed, zerofill, precision=1},
      columns/triad_total/.style={column name={triad incl.\ WA}, fixed, zerofill, precision=1},
      every head row/.style={before row={\toprule & & \multicolumn{4}{c}{GiB/s}\\}, after row=\midrule},
      every last row/.style={after row=\bottomrule},
    ]{data/ceilings_bw.csv}%
  }
\end{table}

\begin{table}[tbp]
  \centering\small
  \caption{Arithmetic intensity (float32) from the actual traffic per element, the same convention
  as the triad roof: L2, L3 and DRAM include the write-allocate read of a store to an unread line;
  in L1 every store hits, so there is none.}\label{tab:roofline-inputs}
  \begin{tabularx}{\linewidth}{@{}lrrX@{}}
    \toprule
    Kernel & $I$, L2--DRAM & $I$, L1 & Traffic per element (L2--DRAM) \\
    \midrule
    SAXPY        & \num{0.167} & \num{0.167} & $2\,\text{FLOP} / \qty{12}{\byte}$: read $x$, read $y$, write $y$
                                  ($y$ is read first, so no write-allocate) \\
    multiply     & \num{0.063} & \num{0.083} & $1\,\text{FLOP} / \qty{16}{\byte}$: read $x,y$, write $z$ + write-allocate of $z$ \\
    stencil      & \num{0.417} & \num{0.625} & $5\,\text{FLOP} / \qty{12}{\byte}$: read $x$ once (reuse), write $y$ + write-allocate \\
    dot product  & \num{0.25}  & \num{0.25}  & $2\,\text{FLOP} / \qty{8}{\byte}$: read $x,y$, no store \\
    triad (roof) & \num{0.125} & \num{0.167} & $2\,\text{FLOP} / \qty{16}{\byte}$: read $b,c$, write $a$ + write-allocate \\
    \bottomrule
  \end{tabularx}
\end{table}

\begin{figure}[tbp]
  \centering
  \IfDataExists{roofline_points}{%
    \begin{tikzpicture}
      \begin{axis}[acs, xmode=log, ymode=log, xlabel={Arithmetic intensity (FLOP/byte)},
                   ylabel={Performance (GFLOP/s)}, width=0.7\linewidth, height=0.6\linewidth,
                   xmin=0.02, xmax=20, ymin=0.5, ymax=400,
                   legend pos=outer north east, legend style={font=\scriptsize}]
        \IfFileExists{data/roofline_roofs.csv}{%
          \addplot[thick, black] table[x=ai, y=gflops, only if={roof}{L1}] {data/roofline_roofs.csv};
          \addlegendentry{L1 roof}
          \addplot[thick, black!70, densely dashed] table[x=ai, y=gflops, only if={roof}{L2}] {data/roofline_roofs.csv};
          \addlegendentry{L2 roof}
          \addplot[thick, black!55, dashdotted] table[x=ai, y=gflops, only if={roof}{L3}] {data/roofline_roofs.csv};
          \addlegendentry{L3 roof}
          \addplot[thick, black!40, dotted] table[x=ai, y=gflops, only if={roof}{DRAM}] {data/roofline_roofs.csv};
          \addlegendentry{DRAM roof}
          \addplot[thin, red!60] table[x=ai, y=gflops, only if={roof}{peak_sse_f32}] {data/roofline_roofs.csv};
          \addlegendentry{SSE peak (autovec)}
          \addplot[thin, red!60, densely dashed] table[x=ai, y=gflops, only if={roof}{peak_scalar_f32}] {data/roofline_roofs.csv};
          \addlegendentry{scalar peak}
        }{}
        \addplot+[only marks, mark=*, nodes near coords, point meta=explicit symbolic,
                  every node near coord/.style={font=\tiny, anchor=south west}, iqr bars]
          table[x=ai, y=gflops_median, iqr of=gflops, meta=kernel, only if={variant}{native}, only if={level}{L1}] {data/roofline_points.csv};
        \addlegendentry{native, L1-resident}
        \addplot+[only marks, mark=square*, iqr bars] table[x=ai, y=gflops_median, iqr of=gflops, only if={variant}{native}, only if={level}{L2}] {data/roofline_points.csv};
        \addlegendentry{native, L2-resident}
        \addplot+[only marks, mark=triangle*, iqr bars] table[x=ai, y=gflops_median, iqr of=gflops, only if={variant}{native}, only if={level}{L3}] {data/roofline_points.csv};
        \addlegendentry{native, L3-resident}
        \addplot+[only marks, mark=diamond*, iqr bars] table[x=ai, y=gflops_median, iqr of=gflops, only if={variant}{native}, only if={level}{DRAM}] {data/roofline_points.csv};
        \addlegendentry{native, DRAM}
        \addplot+[only marks, mark=o, black, iqr bars] table[x=ai, y=gflops_median, iqr of=gflops, only if={variant}{scalar}, only if={level}{DRAM}] {data/roofline_points.csv};
        \addlegendentry{scalar, DRAM}
      \end{axis}
    \end{tikzpicture}%
  }
  \caption{Cache-aware roofline for one core (L1 points at their L1 intensity,
  \cref{tab:roofline-inputs}). Roofs: measured triad bandwidth per level, capped by the
  measured AVX2-FMA float32 peak; thin red lines: the measured SSE and scalar peaks that bound the
  \code{autovec} and \code{scalar} builds. Points: kernel GFLOP/s from \cref{sec:baseline} at
  $N=\num{1024}$, \num{16384}, \num{524288} and \num{16777216} (L1, L2, L3, DRAM), median with IQR error bars (mostly smaller than the markers).}\label{fig:roofline}
\end{figure}

\subsection{Analysis}
\emph{Ceilings} (\cref{tab:ceilings}). The AVX2 FMA loop reaches \qty{89.2}{\giga\flop\per\second}
float32 (\num{44.6} float64), which is \num{31.99}~FLOP/cycle at the measured
\qty{2.79}{\giga\hertz}. That is exactly the $2~\text{FMA ports} \times 8~\text{lanes} \times 2 = 32$
FLOP/cycle the core can do, so the ports, not the loop, set this ceiling. The SSE mul+add loop
(\code{autovec}'s ISA, no FMA) reaches \qty{22.3}{\giga\flop\per\second} (\num{8.00}~FLOP/cycle,
expected $2\times4=8$) and the scalar loop \qty{5.55}{\giga\flop\per\second}
(\num{1.99}~FLOP/cycle, expected $2$). Triad bandwidth (actual traffic) falls from
\qty{155.6}{\gibi\byte\per\second} in L1 to \num{90.5}, \num{48.6} and
\qty{24.3}{\gibi\byte\per\second} in L2, L3 and DRAM. In L1 this is the explicit load and store
traffic: the triad's stores hit lines already in L1, so the write-allocate read counted below L1
(\qty{207}{\gibi\byte\per\second} incl.\ WA) does not happen there. The ridge points are
$I = P_{\text{peak}}/B_\ell \approx \num{0.53}$, \num{0.92}, \num{1.71} and
\qty{3.42}{\flop\per\byte}. The L1 ridge is just above the \qty{0.5}{\flop\per\byte} the second
hypothesis assumed. Only the stencil, whose L1 traffic is \qty{8}{\byte} per element
($I=\num{0.625}$), lies right of it, i.e.\ nominally under the compute roof in L1.

\emph{DRAM-resident kernels.} Every kernel's intensity (\numrange{0.06}{0.42}~FLOP/byte) is far left
of the DRAM ridge, so all four are bandwidth-bound in DRAM. This is consistent with the hypothesis
and with the locality sweep (\cref{sec:locality}). Against the DRAM roof $B_{\text{DRAM}}\cdot I$,
native multiply reaches \SI{100}{\percent} (\num{1.63} measured vs \qty{1.63}{\giga\flop\per\second}),
the stencil \SI{101}{\percent} and SAXPY \SI{107}{\percent} (\num{4.65} vs \num{4.35}). Multiply has
exactly the triad's access pattern (two reads, one write to an unread line) and lands on the roof.
SAXPY, whose pattern differs (two reads, one write to a line it just read), exceeds it by
\SI{7}{\percent}, and by up to \SI{20}{\percent} at the intermediate $N=\num{8388608}$
(\cref{sec:negative}). The sustainable bandwidth depends a little on the access pattern, so the
triad roof is a reference, not an upper bound to the byte. Because the roof comes from a different
loop, a kernel landing on it is evidence, not a tautology. The scalar SAXPY, multiply and stencil
DRAM points sit at \SIrange{37}{66}{\percent} of the same roof (scalar dot: \SI{20}{\percent},
latency-bound like native dot). They are the largely instruction-bound loops of
\cref{sec:locality}, which is why SIMD still pays $1.51\times$
(multiply), $1.94\times$ (SAXPY) and $2.72\times$ (stencil) at these points.

\emph{Cache-resident kernels.} The L3-resident native points sit at \SIrange{95}{100}{\percent} of the
L3 roof (SAXPY \SI{100}{\percent}, multiply \SI{100}{\percent}, stencil \SI{95}{\percent}), so L3
bandwidth bounds them exactly. Closer to the core they fall below the roof. In L2, SAXPY reaches
\SI{86}{\percent} and multiply \SI{88}{\percent} of the roof (stencil \SI{42}{\percent}). In L1,
native SAXPY delivers \qty{20.1}{\giga\flop\per\second} at $I=\num{0.167}$, \SI{72}{\percent} of the
L1 roof and \SI{23}{\percent} of the FMA peak; multiply reaches \SI{66}{\percent} of the L1 roof and
the stencil only \SI{19}{\percent} of the peak that bounds it there. So cache-resident kernels are
bounded neither by the FMA units (confirming the second hypothesis) nor quite by the measured L1
bandwidth. The limit is the loop's own load/store and instruction issue:
\numrange{0.22}{0.28}~cycles per element for SAXPY versus the \num{0.125} that two loads and a store
per cycle would allow (\cref{sec:baseline}). The roofline thus explains the step pattern of
\cref{fig:locality} from L2 outward, since each step is a move down to a lower roof. Inside L1/L2 it
only gives the upper bound.

\emph{The dot product} is the one kernel far \emph{below} every roof:
\qtyrange{1.33}{1.49}{\giga\flop\per\second} at all sizes (\cref{sec:baseline}). It is not
bandwidth-bound but latency-bound by the non-reassociated reduction (\cref{sec:vectorization}). One
scalar \code{vaddss} per element with a \num{4}-cycle latency gives an in-core ceiling of
$2~\text{FLOP}/4~\text{cycles} = 0.5$~FLOP/cycle, i.e.\ \qty{1.39}{\giga\flop\per\second} at the
measured \qty{2.79}{\giga\hertz}. The measured \qtyrange{1.39}{1.42}{\giga\flop\per\second} for every
$N$ from \num{1024} to \num{1048576} matches it to within \SI{2}{\percent}; only $N\le512$ runs slightly
above it (\qtyrange{1.45}{1.49}{\giga\flop\per\second}). Each call starts a new, independent chain,
and with short chains the out-of-order core overlaps the end of one call with the next. DRAM runs
slightly below it (\qty{1.33}{\giga\flop\per\second}). The roofline, a throughput model, does not apply to the dot
product until the dependency chain is broken. To move a streaming kernel toward the compute roof one
would have to raise its arithmetic intensity (fuse passes, or block for cache so data is re-read from
L1 rather than DRAM); widening the vectors alone cannot help a kernel that is already waiting on
memory.

% !TeX root = ../main.tex
\section{Negative Example}\label{sec:negative}
\rubric{AI-use transparency and failure case}{10}{A meaningful negative or misleading example, verified
and explained (with \cref{sec:ai-use}).}

\begin{hypothesis}
  \todo{What a naive reading (or a quickly generated benchmark) would predict for this case.}
\end{hypothesis}

\begin{knobs}
  Case        & \todo{a loop that looks SIMD-friendly but is not vectorized, or one that vectorizes but gains little because memory bandwidth dominates} \\
  Build       & \todo{variant / extra flags} \\
  Script      & \code{SIMD_Advantage_Profiling/scripts/run.sh negative} \\
\end{knobs}

\subsection{Evidence}
\todo{Vectorization report or disassembly excerpt plus timing figure/table.}

\subsection{Explanation}
\todo{Why the expected speedup does not appear, and how you verified the explanation.}

% !TeX root = ../main.tex
\section{Anomalies and Limitations}\label{sec:anomalies}
\rubric{Report quality}{10}{Clear plots, concise writing, honest anomalies and thoughtful limitations.}

\subsection{Unexpected results}
\begin{itemize}
  \item \textbf{Dot product never speeds up.} We expected a reduction to vectorize at $\approx8\times$
        like the other float32 kernels. Instead native dot stays at $\approx1\times$ across all sizes.
        Hypothesis: without \code{-ffast-math} the compiler may not reassociate floating-point
        addition, so it vectorizes the multiply but keeps a serial scalar sum. Tested: the
        disassembly shows \code{vmulps} on \code{ymm} followed by a chain of scalar \code{vaddss}
        (\cref{sec:vectorization}), confirming the hypothesis. We chose not to enable
        \code{-ffast-math} so that all variants remain bit-reproducible against the reference.
  \item \textbf{Native SAXPY speedup varies within L1.} It is $5.6\text{--}5.9\times$ at
        $N=\num{1024}$, \num{2048}, \num{3072} and \num{4096}, but $6.86\times$ at $N=\num{1536}$
        (locality sweep) and $6.32\times$ at $N=512$ (baseline). Scalar time is flat across these
        sizes and L1D misses stay below \num{0.0011} per element below $N=\num{4096}$; the peak is
        native needing \num{0.22} instead of \numrange{0.26}{0.28}~cycles per element. Per-call
        overhead cannot explain it, since it \emph{lowers} small-$N$ speedups
        (\cref{sec:alignment}). Untested hypothesis: the relative placement of \code{x} and
        \code{y} (4K aliasing, as for the stencil below), testable with the same distance sweep.
  \item \textbf{Stencil keeps a $2.6\text{--}2.8\times$ advantage in DRAM} while SAXPY and
        multiply fall to $1.88\text{--}2.02\times$ and $1.5\text{--}1.6\times$. Its native build is
        just as bandwidth-bound as theirs: it runs at \SI{101}{\percent} of the DRAM roof
        (\cref{sec:roofline}). Its scalar build, however, needs
        \num{5}~FLOP per element of serial \code{mulss}/\code{addss} work and reaches only
        \qtyrange{4.0}{5.0}{\giga\flop\per\second} at every size (\qty{4.0}{\giga\flop\per\second},
        \SI{37}{\percent} of the roof, in DRAM). The retained
        ratio is again scalar's instruction cost, as for SAXPY (\cref{sec:locality}), just larger.
  \item \textbf{Cache-resident stencil is slower than L3-resident stencil.} Native stencil delivers
        \qtyrange{16.6}{17.3}{\giga\flop\per\second} for $N=\num{1024}\ldots\num{16384}$ but
        \qtyrange{19.9}{20.7}{\giga\flop\per\second} for $N=\num{65536}\ldots\num{1048576}$, and
        \num{31.4} and \qty{27.4}{\giga\flop\per\second} at the isolated sizes $N=512$ and
        \num{32768}. The pattern reproduces exactly across separate runs; scalar stencil stays at \qtyrange{4.0}{5.0}{\giga\flop\per\second}. A cache
        effect would make the rate fall monotonically with $N$; this does not. Hypothesis (untested): 4K aliasing. The small \code{x} and \code{y}
        allocations sit at addresses whose difference makes a store to \code{y[i]} look like it
        conflicts with the load of \code{x[i+1]} in the low 12 address bits, which stalls the load;
        large allocations come from fresh pages with a different relative offset. It would be confirmed by a
        sweep of the \code{x}--\code{y} distance modulo \qty{4}{\kibi\byte} at fixed $N$, showing
        the slow rate only at small distances, together with the core's 4K-alias counter
        (\code{ld_blocks_partial.address_alias}). We ran neither.
  \item \textbf{Scalar and SIMD SAXPY do not converge in DRAM.} We expected one shared bandwidth
        plateau; native stays $1.88\text{--}2.02\times$ ahead because scalar SAXPY is still largely
        instruction-bound there (\num{2.34}~cycles per element at IPC \num{3.0}) while native sits on
        the roof (IPC \num{0.6}); the cycle ratio, $1.88$, is the measured speedup
        (\cref{sec:locality}).
  \item \textbf{Same code, different process, $N=\num{32768}$.} Scalar SAXPY is stable across
        processes elsewhere (\qtyrange{0.650}{0.653}{\nano\second} per element at all \num{9}
        locality sizes from $N=\num{65536}$ to \num{1048576}), but where the two arrays exactly fill
        the \qty{256}{\kibi\byte} L2 two measurements of the same machine code differ by
        \SI{1}{\percent} (scalar) and \SI{14}{\percent} (native); page placement is the untested
        hypothesis (\cref{sec:stride}).
  \item \textbf{Hardware gather loses to scalar.} The explicit \code{vgatherdps} kernel takes
        \qty{3.23}{\nano\second} per element at $N=\num{32768}$ versus \num{2.16} (native scalar
        loop) and \qty{2.02}{\nano\second} (scalar build), i.e.\ $1.5\text{--}1.6\times$
        \emph{slower}, and it is \SI{7}{\percent} slower in DRAM. Hypothesis: the Gather Data
        Sampling microcode mitigation, which this CPU runs (\cref{sec:stride}) and which is known
        to make \code{vgatherdps} much slower, plus the eight scalar
        stores that AVX2 cannot vectorize. Confirming it would need a boot without the mitigation,
        which requires root.
\end{itemize}

\subsection{Limitations}
\begin{itemize}
  \item \textbf{Core-PMU counters only.} The SIMD locality counters are core-PMU, user-space
        events of the benchmark thread (\cref{sec:ceilings-method}): cycles, instructions, generic LLC
        misses and L1D read misses. We did not count DRAM traffic for SIMD. The machine can do it:
        the companion Cache \& Memory report, run on the same machine the same day with
        \code{perf_event_paranoid}~$=0$, reads the memory-controller (\code{uncore_imc}) counters
        and finds IMC traffic matching its streaming estimate to within about \SI{1}{\percent}. TLB
        events were not recorded either. The generic LLC-miss event's treatment of hardware-prefetch
        requests is model-specific, so we use misses per element as a relative measure between
        builds rather than as exact DRAM traffic. Counting IMC bytes per element for scalar and
        native SAXPY would turn the roofline's traffic assumption into a measurement.
  \item \textbf{Fixed clock.} All reported experiments were run on 2026-10-21 with governor
        \sysGovernor{}, turbo \sysTurbo{} and SMT off; the cycle counter confirms
        \qtyrange{2.78}{2.79}{\giga\hertz} for every locality and ceilings configuration
        (\cref{sec:method}). The fixed clock makes the results reproducible, but it is not how the
        CPU normally runs. With turbo the compute-bound numbers scale up by roughly the clock ratio
        ($1.72\text{--}1.75\times$ in our turbo sweep) while DRAM bandwidth does not, so the
        DRAM-resident SAXPY speedups reported here ($1.88\text{--}2.02\times$) are larger than at turbo
        ($1.32\text{--}1.44\times$, \cref{sec:method}).
  \item \textbf{One machine, one compiler.} All results are from a single i9-10900F with GCC
        \sysGCC{}. Clang was used only to confirm the build compiles. No AVX-512 part was available,
        so the \qty{512}{bit} regime is out of scope.
  \item \textbf{THP always-on; no core isolation.} Transparent huge pages are \code{always}, and
        the pinned core is not isolated from the scheduler (no \code{isolcpus}). Neither is controlled
        for; the lock file only serializes cooperating jobs, so the shared-core checks of
        \cref{sec:method} (repetition spread $\le1.36\times$, apparent clock
        \qtyrange{2.78}{2.79}{\giga\hertz}) are what guard the reported data.
  \item \textbf{Discarded sessions.} Earlier sessions the same day were discarded: those before the
        GPE storm was masked (\cref{sec:method}), and a first clean session during which another
        agent's unlocked smoke tests briefly shared CPU~2 ($\approx2\times$ slowdowns in some
        repetitions).
  \item \textbf{Report build backend.} Without \code{biber}, \code{biblatex} uses the BibTeX
        backend, so the single reference is sorted by name/year to avoid a spurious rerun warning.
\end{itemize}

% !TeX root = ../main.tex
\section{AI-Use Note}\label{sec:ai-use}

Summarized from \code{AI_USE_LOG.md} in the code repository.

\begin{description}[style=nextline]
  \item[Tools used] \todo{Tools and models, and for which parts of the project.}
  \item[What they helped generate] \todo{e.g.\ repository scaffolding, build and run scripts, debugging help.}
  \item[What was independently modified, tested, or rejected] \todo{Concrete examples, including at
        least one AI suggestion that was wrong or misleading and how it was caught.}
  \item[Validation] \todo{How correctness and measurements were checked independently (reference
        comparison, vectorization evidence, sanity checks).}
\end{description}

% !TeX root = ../main.tex
\section{Conclusion}\label{sec:conclusion}

\todo{Three to five sentences: when SIMD helps on this machine, when gains compress or vanish, and
the single most important lesson for writing and benchmarking vectorized code.}

% % !TeX root = ../main.tex
% Optional CPU-vs-GPU extension. Enable with % !TeX root = ../main.tex
% Optional CPU-vs-GPU extension. Enable with % !TeX root = ../main.tex
% Optional CPU-vs-GPU extension. Enable with \input{sections/15_gpu_extension} in main.tex.
\section{Optional Extension: CPU vs.\ GPU}\label{sec:gpu}

GPU: \sysGPU; CUDA \sysCUDA.

\begin{hypothesis}
  \todo{Break-even $N$ you expect, once host--device transfer is included.}
\end{hypothesis}

\begin{knobs}
  Kernel      & \tbd \\
  Timing      & kernel-only (device synchronized) and end-to-end (including host$\leftrightarrow$device copies) \\
  Varied      & $N$ \\
\end{knobs}

\todo{Figure: CPU (best SIMD build) vs.\ GPU kernel-only vs.\ GPU end-to-end throughput across $N$.}

\subsection{Analysis}
\todo{When GPU acceleration helps, when transfers erase the benefit, and why.}
 in main.tex.
\section{Optional Extension: CPU vs.\ GPU}\label{sec:gpu}

GPU: \sysGPU; CUDA \sysCUDA.

\begin{hypothesis}
  \todo{Break-even $N$ you expect, once host--device transfer is included.}
\end{hypothesis}

\begin{knobs}
  Kernel      & \tbd \\
  Timing      & kernel-only (device synchronized) and end-to-end (including host$\leftrightarrow$device copies) \\
  Varied      & $N$ \\
\end{knobs}

\todo{Figure: CPU (best SIMD build) vs.\ GPU kernel-only vs.\ GPU end-to-end throughput across $N$.}

\subsection{Analysis}
\todo{When GPU acceleration helps, when transfers erase the benefit, and why.}
 in main.tex.
\section{Optional Extension: CPU vs.\ GPU}\label{sec:gpu}

GPU: \sysGPU; CUDA \sysCUDA.

\begin{hypothesis}
  \todo{Break-even $N$ you expect, once host--device transfer is included.}
\end{hypothesis}

\begin{knobs}
  Kernel      & \tbd \\
  Timing      & kernel-only (device synchronized) and end-to-end (including host$\leftrightarrow$device copies) \\
  Varied      & $N$ \\
\end{knobs}

\todo{Figure: CPU (best SIMD build) vs.\ GPU kernel-only vs.\ GPU end-to-end throughput across $N$.}

\subsection{Analysis}
\todo{When GPU acceleration helps, when transfers erase the benefit, and why.}
   % optional CPU-vs-GPU comparison

\printbibliography

\appendix
% !TeX root = ../main.tex
\section{Reproducibility}\label{app:repro}

Code, scripts and raw data: \coderepo, on the \code{class-wide-project} submission branch
(the exact commit is recorded in \code{git log}).

\begin{lstlisting}[language=bash]
git clone https://github.com/SKempiak/ACS-Code.git && cd ACS-Code
git checkout <commit>
make -C SIMD_Advantage_Profiling          # scalar / autovec / native builds
make tools                                # tools/summarize
make -C SIMD_Advantage_Profiling test     # 570 correctness cases
make -C SIMD_Advantage_Profiling report   # vectorization evidence
scripts/collect_env.sh --tex > ../ACS-Latex/common/system_info.tex
SIMD_Advantage_Profiling/scripts/run.sh all
scripts/sync_data.sh SIMD_Advantage_Profiling
\end{lstlisting}

\begin{table}[h]
  \centering\small
  \caption{Raw data behind each figure (directories under
  \code{SIMD_Advantage_Profiling/data/raw/}).}\label{tab:raw-data}
  \begin{tabularx}{\linewidth}{@{}lX@{}}
    \toprule
    Figure / table & Raw run directory \\
    \midrule
    \Cref{fig:baseline-speedup,fig:baseline-throughput} & \code{baseline/20261012-182813} \\
    \Cref{fig:locality} & \code{locality/20261012-182813} \\
    \Cref{fig:alignment} & \code{alignment/20261012-182813} \\
    \Cref{fig:stride}, \cref{tab:stride} & \code{stride/20261012-182813} \\
    \Cref{fig:dtype} & \code{dtype/20261012-182813} \\
    \Cref{fig:roofline} & \code{baseline/20261012-182813}, \code{negative/20261012-182813} \\
    \Cref{tab:vec-summary} & \code{vectorization/20261012-182734} \\
    \bottomrule
  \end{tabularx}
\end{table}

Each raw directory holds \code{raw.csv} (one row per repetition, with full configuration metadata)
and \code{env.txt} (a snapshot of \code{collect_env.sh}); the processed CSVs plotted above are
produced from them by \code{tools/summarize} and copied into this report's \code{data/} by
\code{scripts/sync_data.sh}. Every number quoted in the report traces to one of these files.


\end{document}
